\scanpage{28}

techniquement simple. Si $X$ est un \deleted{\unsure{paracompact}} espace paracompact, on pourra prendre (à titre provisoire !) pour $\unsure{\mathcal V_{Y,X}|V_{Y,X}}$ le ``germe d'espaces'' de $X$ autour de $Y$ -- qui est donc un topos, sav. la $\varprojlim$ toposique des voisinages ouverts de $Y$ dans $X$. Intuitivement, $\unsure{\mathcal V_{Y,X}|V_{Y,X}}\hookrightarrow X$ a les propriétés d'une imm. ouverte, et $Y\hookrightarrow\unsure{\mathcal V_{Y,X}|V_{Y,X}}$ celles d'une immersion fermée, et (3) constitue une factorisation canonique de (1).

\deleted{L'hypothèse de locale trivialité normale implique que l'on a (moyennant des conditions de ``modération'' convenables) une inclusion}

\margin{L'intérêt principal de cette factorisation, c'est que \unsure{l'inclusion} $Y\longrightarrow\unsure{\mathcal V_{Y,X}|V_{Y,X}}$ est une ``équivalence d'homotopie'' (indépendamment de toute hyp. de locale trivialité).}

(4)
\diagram{\begin{tikzcd}[ampersand replacement=\&]
Y\arrow[r,"\text{équi. d'homotopie}"]\&\unsure{\mathcal V_{Y,X}|V_{Y,X}}
\end{tikzcd}}

Supposons p. ex. $Y$ et $X-Y=X^*$ des variétés, et $Z\inserted{{}^*=Z\setminus\{z\}}$ une variété (moyennant des hyp. de modération, l'hyp. sur $Z$ devrait être une conséquence de la première hypothèse), \deleted{\unsure{et que}} alors que $z$ est un pt singulier \inserted{de $Z$}, i.e. $X$ est ``singulier le long de $Y$''. Il est naturel alors de vouloir reconstruire $X$ en termes des variétés $Y$ et $X^*$,

\margin{27}
